\documentclass{beamer}

\mode<presentation>
{
  \usetheme{Madrid}
  \usecolortheme{rose}
  \usefonttheme{structurebold}
  \setbeamertemplate{navigation symbols}{}
  \setbeamertemplate{caption}[numbered]
}

\usepackage[vietnamese]{babel}
\usepackage{fontspec}
\usepackage{tikz}
\usepackage{xcolor}
\usetikzlibrary{arrows.meta,calc,positioning}

\setmainfont{Arial}
\setsansfont{Arial}
\setmonofont{Consolas}
\newfontfamily\displayfont{Times New Roman}

\definecolor{AppleBlue}{HTML}{0071E3}
\definecolor{AppleGreen}{HTML}{34C759}
\definecolor{AppleOrange}{HTML}{FF9500}
\definecolor{AppleMuted}{HTML}{6E6E73}
\definecolor{AppleInk}{HTML}{1D1D1F}
\definecolor{AppleBorder}{HTML}{D2D2D7}
\definecolor{AppleSurface}{HTML}{F5F5F7}

\setbeamercolor{structure}{fg=AppleBlue}
\setbeamercolor{normal text}{fg=AppleInk,bg=white}
\setbeamercolor{block title}{fg=white,bg=AppleBlue}
\setbeamercolor{block body}{fg=AppleInk,bg=AppleSurface}
\setbeamercolor{exampleblock title}{fg=white,bg=AppleGreen}
\setbeamercolor{exampleblock body}{fg=AppleInk,bg=AppleSurface}
\setbeamercolor{alertblock title}{fg=white,bg=AppleInk}
\setbeamercolor{alertblock body}{fg=AppleInk,bg=AppleSurface}
\setbeamerfont{frametitle}{family=\displayfont,size=\Large}
\setbeamerfont{title}{family=\displayfont,size=\LARGE}

\title[Hilbert \& Qubit Tuần 1]{Hilbert \& Qubit Tuần 1}
\author{Self-learning QML}
\date{\today}

\AtBeginSection[]
{
  \begin{frame}[plain]
    \tableofcontents[currentsection]
  \end{frame}
}

\begin{document}

\begin{frame}
  \titlepage
\end{frame}

\begin{frame}[plain]{Outline}
  \tableofcontents
\end{frame}

\section{Nền tảng}

\begin{frame}{Vì sao cần Hilbert?}
  \begin{columns}[T,onlytextwidth]
    \column{0.48\textwidth}
    \begin{block}{Bit cổ điển}
      Bit chỉ có một trạng thái tại một thời điểm: `0` hoặc `1`.
    \end{block}
    \begin{block}{Qubit}
      Qubit là vector phức đã chuẩn hóa, nên có thể mang chồng chất và pha.
    \end{block}
    \begin{block}{Mục tiêu tuần 1}
      Đọc được qubit bằng 3 ngôn ngữ: đại số, hình học và code.
    \end{block}

    \column{0.48\textwidth}
    \centering
    \begin{tikzpicture}[scale=1.0]
      \draw[AppleBorder, line width=0.8pt] (0,0) circle (1.55);
      \draw[AppleBorder, line width=0.55pt] (0,0) ellipse (1.55 and 0.55);
      \draw[AppleBorder, line width=0.55pt] (0,-1.55) arc (-90:90:0.55 and 1.55);
      \draw[AppleBlue, line width=1.2pt, -stealth] (0,0) -- (0.95,1.1);
      \fill[AppleBlue] (0.95,1.1) circle (2pt);
      \node[above] at (0,1.75) {$|0\rangle$};
      \node[below] at (0,-1.75) {$|1\rangle$};
      \node[right] at (1.15,1.1) {$|\psi\rangle$};
    \end{tikzpicture}
  \end{columns}
\end{frame}

\begin{frame}{Trạng thái qubit}
  \begin{columns}[T,onlytextwidth]
    \column{0.46\textwidth}
    \begin{block}{Dạng tổng quát}
      \small
      \[
        \begin{aligned}
          |\psi\rangle &= \alpha|0\rangle + \beta|1\rangle \\
          |\alpha|^2 + |\beta|^2 &= 1
        \end{aligned}
      \]
    \end{block}
    \begin{exampleblock}{Ý chính}
      Hệ số là số phức, còn bình phương biên độ mới là xác suất.
    \end{exampleblock}

    \column{0.48\textwidth}
    \begin{tikzpicture}[x=1cm,y=1cm,scale=0.86,transform shape]
      \draw[draw=AppleBorder, fill=AppleSurface, rounded corners=8pt] (0,0) rectangle (5.0,2.6);
      \draw[fill=AppleBlue, draw=AppleBlue] (0.85,0.45) rectangle (1.55,1.85);
      \draw[fill=AppleGreen, draw=AppleGreen] (2.2,0.45) rectangle (3.35,1.15);
      \node[anchor=south] at (1.2,1.85) {\scriptsize $|\alpha|^2$};
      \node[anchor=south] at (2.8,1.15) {\scriptsize $|\beta|^2$};
      \node at (1.2,0.2) {\scriptsize $|0\rangle$};
      \node at (2.8,0.2) {\scriptsize $|1\rangle$};
      \node[anchor=west] at (3.8,1.5) {\small vector trạng thái};
      \node[anchor=west] at (3.8,1.1) {\small dự đoán đo};
      \node[anchor=west] at (3.8,0.7) {\small trước khi đo};
    \end{tikzpicture}
  \end{columns}
\end{frame}

\begin{frame}{Bra-ket}
  \begin{columns}[T,onlytextwidth]
    \column{0.31\textwidth}
    \begin{block}{Ket}
      \centering {\Large $|\psi\rangle$}\par
      Vector cột.
    \end{block}
    \column{0.31\textwidth}
    \begin{block}{Bra}
      \centering {\Large $\langle\psi|$}\par
      Liên hợp chuyển vị.
    \end{block}
    \column{0.31\textwidth}
    \begin{block}{Outer product}
      \centering {\Large $|\psi\rangle\langle\phi|$}\par
      Tạo toán tử hoặc $\rho$.
    \end{block}
  \end{columns}

  \vspace{0.5em}
  \begin{center}
    \begin{tikzpicture}[node distance=1.4cm]
      \node[draw=AppleBorder, fill=AppleSurface, rounded corners=8pt, minimum width=3cm, minimum height=1cm] (k) {$|\psi\rangle$};
      \node[draw=AppleBorder, fill=AppleSurface, rounded corners=8pt, right=2.0cm of k, minimum width=3cm, minimum height=1cm] (b) {$\langle\psi|$};
      \draw[-{Stealth[length=2.2mm]}, line width=0.8pt, AppleBlue] (k) -- node[above, font=\scriptsize] {liên hợp chuyển vị} (b);
    \end{tikzpicture}
  \end{center}
\vspace{0.15em}
\begin{block}{Công thức lõi}
  \small
  \[
    \begin{aligned}
      |\psi\rangle &= \alpha|0\rangle + \beta|1\rangle \\
      \langle\psi|\psi\rangle &= 1
    \end{aligned}
  \]
\end{block}
\end{frame}

\section{Bloch}

\begin{frame}{Mặt cầu Bloch}
  \begin{columns}[c,onlytextwidth]
    \column{0.46\textwidth}
    \centering
    \begin{tikzpicture}[scale=1.08]
      \draw[AppleBorder, line width=0.8pt] (0,0) circle (2);
      \draw[AppleBorder, line width=0.55pt] (0,0) ellipse (2 and 0.7);
      \draw[AppleBorder, line width=0.55pt] (0,-2) arc (-90:90:0.7 and 2);
      \draw[AppleBlue, line width=1.2pt, -stealth] (0,0) -- (1.2,1.35);
      \fill[AppleBlue] (1.2,1.35) circle (2pt);
      \node[above] at (0,2.15) {$|0\rangle$};
      \node[below] at (0,-2.2) {$|1\rangle$};
      \node[right] at (2.1,0) {$x$};
      \node[anchor=west] at (1.3,1.45) {$\theta,\phi$};
    \end{tikzpicture}

    \column{0.48\textwidth}
    \begin{block}{Hai góc đủ để định vị}
      \begin{itemize}
        \item $\theta$ điều khiển vị trí theo trục $z$
        \item $\phi$ là pha tương đối
      \end{itemize}
    \end{block}
    \begin{exampleblock}{Một ví dụ}
      Nếu $|\alpha|=|\beta|$, trạng thái nằm trên xích đạo.
    \end{exampleblock}
  \end{columns}
\end{frame}

\begin{frame}{Pha tương đối}
  \begin{columns}[T,onlytextwidth]
    \column{0.44\textwidth}
    \begin{block}{Cùng xác suất}
      \centering
      \begin{tikzpicture}[scale=0.85]
        \draw[AppleBorder, line width=0.7pt] (0,0) circle (1.1);
        \draw[AppleBlue, line width=1.0pt, -stealth] (0,0) -- (0.85,0.55);
      \end{tikzpicture}
    \end{block}
    \column{0.44\textwidth}
    \begin{block}{Khác pha}
      \centering
      \begin{tikzpicture}[scale=0.85]
        \draw[AppleBorder, line width=0.7pt] (0,0) circle (1.1);
        \draw[AppleOrange, line width=1.0pt, -stealth] (0,0) -- (0.15,0.99);
        \draw[AppleGreen, line width=1.0pt, -stealth] (0,0) -- (-0.55,0.85);
      \end{tikzpicture}
    \end{block}
  \end{columns}

  \begin{exampleblock}{Chốt}
    Pha khác nhau làm giao thoa khác nhau, dù xác suất đo cơ bản giống nhau.
  \end{exampleblock}
\end{frame}

\section{Tiến hóa}

\begin{frame}{Unitary}
  \begin{theorem}
    Một toán tử unitary thỏa mãn
    \[
      U^\dagger U = I
    \]
    và bảo toàn chuẩn, tích vô hướng, cũng như khả năng đảo ngược của tiến hóa.
  \end{theorem}

  \vspace{0.4em}
  \begin{center}
    \begin{tikzpicture}[node distance=1.1cm]
      \node[draw=AppleInk, fill=AppleInk, text=white, rounded corners=10pt, minimum width=3cm, minimum height=1.15cm] (u) {$U^\dagger U = I$};
      \node[draw=AppleBorder, fill=AppleSurface, rounded corners=8pt, font=\scriptsize, below left=0.65cm and 1.5cm of u, minimum width=2.8cm, minimum height=0.9cm] (a) {bảo toàn chuẩn};
      \node[draw=AppleBorder, fill=AppleSurface, rounded corners=8pt, font=\scriptsize, below=0.65cm of u, minimum width=2.8cm, minimum height=0.9cm] (b) {bảo toàn góc};
      \node[draw=AppleBorder, fill=AppleSurface, rounded corners=8pt, font=\scriptsize, below right=0.65cm and 1.5cm of u, minimum width=2.8cm, minimum height=0.9cm] (c) {đảo ngược được};
      \draw[-{Stealth[length=2.2mm]}, AppleBlue] (u) -- (a);
      \draw[-{Stealth[length=2.2mm]}, AppleBlue] (u) -- (b);
      \draw[-{Stealth[length=2.2mm]}, AppleBlue] (u) -- (c);
    \end{tikzpicture}
  \end{center}
\end{frame}

\begin{frame}{Hadamard}
  \begin{columns}[T,onlytextwidth]
    \column{0.40\textwidth}
    \begin{block}{Ma trận}
      \[
        H = \frac{1}{\sqrt{2}}
        \begin{bmatrix}
          1 & 1 \\
          1 & -1
        \end{bmatrix}
      \]
    \end{block}
    \column{0.52\textwidth}
    \begin{block}{Hiệu ứng}
      \[
        H|0\rangle = \frac{|0\rangle + |1\rangle}{\sqrt{2}}
      \]
    \end{block}
    \begin{exampleblock}{Ý nghĩa}
      Hadamard đổi cơ sở và tạo chồng chất cân bằng.
    \end{exampleblock}
  \end{columns}
\end{frame}

\section{Nhiều qubit}

\begin{frame}{Tensor product}
  \begin{columns}[T,onlytextwidth]
    \column{0.44\textwidth}
    \begin{block}{Ghép trạng thái}
      \centering
      \begin{tikzpicture}[scale=0.9]
        \node[draw=AppleBorder, fill=AppleSurface, rounded corners=8pt, minimum width=1.2cm, minimum height=1cm] at (0,0.9) {$|0\rangle$};
        \node[draw=AppleBorder, fill=AppleSurface, rounded corners=8pt, minimum width=1.2cm, minimum height=1cm] at (0,-0.1) {$|1\rangle$};
        \node[draw=AppleBorder, fill=AppleSurface, rounded corners=8pt, minimum width=1.2cm, minimum height=1cm] at (2.2,0.9) {$|0\rangle$};
        \node[draw=AppleBorder, fill=AppleSurface, rounded corners=8pt, minimum width=1.2cm, minimum height=1cm] at (2.2,-0.1) {$|1\rangle$};
        \draw[-{Stealth[length=2mm]}, AppleBlue] (1.35,0.4) -- (1.8,0.4);
      \end{tikzpicture}
      \[
        |0\rangle \otimes |1\rangle \in \mathbb{C}^4
      \]
    \end{block}
    \column{0.44\textwidth}
    \begin{block}{Ghép toán tử}
      \centering
      \begin{tikzpicture}[scale=0.9]
        \draw[AppleBorder, fill=AppleSurface, rounded corners=8pt] (0,0) rectangle (3.8,2.0);
        \draw[AppleBlue, line width=0.8pt] (0.65,1.4) -- (3.15,1.4);
        \draw[AppleBlue, line width=0.8pt] (0.65,0.65) -- (3.15,0.65);
        \draw[AppleBorder, line width=0.8pt] (1.9,0.35) -- (1.9,1.75);
      \end{tikzpicture}
      \[
        X \otimes Z
      \]
    \end{block}
  \end{columns}

  \begin{exampleblock}{Chốt}
    Tensor ghép hệ bằng tích, nên số chiều tăng rất nhanh.
  \end{exampleblock}
\end{frame}

\begin{frame}{Bùng nổ chiều}
  \begin{center}
    \begin{tikzpicture}[x=1.2cm,y=0.95cm]
      \foreach \n/\h/\c in {1/1.0/AppleBlue,2/2.0/AppleGreen,3/3.0/AppleOrange,4/4.0/AppleInk}{
        \draw[draw=AppleBorder, fill=\c!12, rounded corners=6pt] (\n-0.35,0) rectangle (\n+0.35,\h);
        \node[below] at (\n,0) {\scriptsize \n};
      }
      \node[anchor=west] at (4.9,2.2) {\Large $2^n$};
      \draw[AppleBlue, line width=1.0pt, -{Stealth[length=2.5mm]}] (4.2,2.2) -- (4.75,2.2);
      \node[anchor=west] at (5.45,2.2) {\small chiều không gian của $n$ qubit};
    \end{tikzpicture}
  \end{center}

  \vspace{0.6em}
  \begin{center}
    \large Một qubit thêm vào là nhân đôi toàn bộ không gian trạng thái.
  \end{center}
\end{frame}

\section{Đo và mật độ}

\begin{frame}{Đo}
  \begin{columns}[T,onlytextwidth]
    \column{0.46\textwidth}
    \begin{block}{Trước đo}
      \centering
      \begin{tikzpicture}[scale=0.9]
        \draw[AppleBorder, fill=AppleSurface, rounded corners=10pt] (0,0) rectangle (4.1,2.2);
        \draw[AppleBlue, line width=1.0pt, -stealth] (1.0,1.0) -- (1.9,1.45);
        \draw[AppleOrange, line width=1.0pt, -stealth] (1.0,1.0) -- (2.0,0.55);
      \end{tikzpicture}
    \end{block}
    \column{0.46\textwidth}
    \begin{block}{Sau đo}
      \centering
      \begin{tikzpicture}[scale=0.9]
        \draw[AppleBorder, fill=AppleSurface, rounded corners=10pt] (0,0) rectangle (4.1,2.2);
        \node[draw=AppleInk, fill=white, rounded corners=8pt, minimum width=1.0cm, minimum height=0.7cm] at (1.4,1.1) {0};
        \draw[AppleBorder, line width=0.9pt, -{Stealth[length=2.2mm]}] (2.2,1.1) -- (3.2,1.1);
        \node[draw=AppleBlue, fill=white, rounded corners=8pt, minimum width=1.0cm, minimum height=0.7cm] at (3.5,1.1) {$|0\rangle$};
      \end{tikzpicture}
    \end{block}
  \end{columns}

  \begin{alertblock}{Quy tắc Born}
    $P(0)=|\alpha|^2,\quad P(1)=|\beta|^2$
  \end{alertblock}
\end{frame}

\begin{frame}{Ma trận mật độ}
  \begin{columns}[T,onlytextwidth]
    \column{0.46\textwidth}
    \begin{block}{Trạng thái thuần}
      \[
        \rho = |\psi\rangle\langle\psi|
      \]
      \centering
      \begin{tikzpicture}[scale=0.9]
        \draw[AppleBorder, fill=AppleSurface, rounded corners=10pt] (0,0) rectangle (4.0,2.3);
        \draw[AppleBlue, line width=1.0pt] (0.95,0.6) rectangle (3.1,1.7);
      \end{tikzpicture}
    \end{block}
    \column{0.46\textwidth}
    \begin{block}{Trạng thái hỗn hợp}
      \[
        \rho = \sum_i p_i |\psi_i\rangle\langle\psi_i|
      \]
      \centering
      \begin{tikzpicture}[scale=0.9]
        \draw[AppleBorder, fill=AppleSurface, rounded corners=10pt] (0,0) rectangle (4.0,2.3);
        \draw[AppleOrange, line width=1.0pt] (0.9,0.55) rectangle (1.8,1.15);
        \draw[AppleGreen, line width=1.0pt] (2.05,0.55) rectangle (3.1,1.45);
        \draw[AppleBlue, line width=1.0pt] (1.45,1.1) rectangle (2.55,1.85);
      \end{tikzpicture}
    \end{block}
  \end{columns}

  \vspace{0.2em}
  \begin{block}{Ba điều phải kiểm}
    \centering $\rho=\rho^\dagger \qquad \mathrm{Tr}(\rho)=1 \qquad \mathrm{Tr}(\rho^2)\le 1$
  \end{block}
\end{frame}

\section{Thực hành}

\begin{frame}[fragile]{NumPy kiểm chứng}
  \begin{columns}[T,onlytextwidth]
    \column{0.46\textwidth}
    \begin{block}{Mã tối thiểu}
      \begingroup\scriptsize
\begin{verbatim}
import numpy as np
X = np.array([
    [0,1],
    [1,0]
], complex)
Z = np.array([
    [1,0],
    [0,-1]
], complex)
XZ = np.kron(X, Z)
print(XZ)
\end{verbatim}
      \endgroup
    \end{block}
    \column{0.46\textwidth}
    \begin{block}{Khi dùng}
      \begin{itemize}
        \item xác nhận tensor product
        \item debug mạch nhỏ
        \item nối công thức với số liệu
      \end{itemize}
    \end{block}
    \begin{exampleblock}{Ý nghĩa}
      Kiểm tra công thức bằng dữ liệu thật thay vì chỉ tin vào trực giác.
    \end{exampleblock}
  \end{columns}
\end{frame}

\begin{frame}{Tự kiểm}
  \begin{columns}[T,onlytextwidth]
    \column{0.31\textwidth}
    \begin{block}{1}
      Vector phức và chuẩn hóa.
    \end{block}
    \column{0.31\textwidth}
    \begin{block}{2}
      $\theta$ và $\phi$.
    \end{block}
    \column{0.31\textwidth}
    \begin{block}{3}
      Kiểm chứng bằng NumPy.
    \end{block}
  \end{columns}

  \vspace{0.35em}
  \begin{exampleblock}{Takeaway}
    Hilbert cho ngôn ngữ, Bloch cho hình dung, unitary cho tiến hóa, tensor cho ghép hệ, đo cho kết quả.
  \end{exampleblock}
\end{frame}

\begin{frame}{Kết luận}
  \begin{center}
    \vfill
    {\displayfont\fontsize{26}{30}\selectfont Hilbert. Bloch. Unitary. Tensor. Đo.\par}
    \vspace{0.75em}
    {\normalsize\color{AppleMuted}Tuần 1 chỉ cần dựng đúng khung này, rồi các tuần sau mới đẩy sang mạch và thuật toán.\par}
    \vfill
  \end{center}
\end{frame}

\end{document}
